% GENERATED FILE. Edit canonical lessons, then run npm run generate:books.
% canonical-source-hash: fnv1a64:2f796f59986ace70
% canonical-lessons: TA-C391-porulataram, TA-C391-maniyam, TA-C391-parrakkurai, TA-C391-viniyoki, TA-C391-nirnayi

\chapter{Public Money}
\label{ch:ta-public-money}
\hlchaptermodality{391}

\begin{chapteropening}
\textbf{By the end of this chapter:} \emph{I can talk about the economy, subsidies and shortages.}

Complete the last lesson of chapter 391: I can talk about the economy, subsidies and shortages.
\end{chapteropening}

\section[\texorpdfstring{\ta{பொருளாதாரம்} (poruḷātāram)}{பொருளாதாரம் (poruḷātāram)}]{\ta{பொருளாதாரம்} (poruḷātāram) — the economy}

\label{lesson:TA-C391-porulataram}



\begin{quote}
Before the new one: say the Tamil for to restructure, then the Tamil for to strengthen.
\end{quote}

\subsection*{You'll want to know: \ta{பொருளாதாரம்}}
\textbf{\ta{பொருளாதாரம்}} — \emph{poruḷātāram} — \textquotedblleft{}the economy\textquotedblright{}.

\begin{quote}
\ta{மழை} \ta{பெய்யாவிட்டால்} \ta{நாட்டின்} \ta{பொருளாதாரம்} \ta{பாதிக்கப்படும்}.
\end{quote}

\begin{quote}
\emph{maḻai peyyāviṭṭāl nāṭṭiṉ poruḷātāram pātikkappaṭum.}
\end{quote}

\begin{quote}
If it does not rain, the country's economy will suffer.
\end{quote}

\ta{பொருளாதார}, economic, before a noun: \ta{பொருளாதார} \ta{வளர்ச்சி}, economic growth.

\begin{cousinweb}
\textbf{\ta{பொருளாதாரம்}} (\emph{poruḷātāram}) is built from native \emph{poruḷ}, \textquotedblleft{}thing, wealth\textquotedblright{} (DEDR 4544), and \emph{ātāram}, \textquotedblleft{}support, basis\textquotedblright{}, borrowed from Sanskrit \emph{ādhāra} (CDIAL 1165): what wealth rests on.
\end{cousinweb}

\begin{grammarlens}[title={What you've built}]
One more word for saying what happens if.
\end{grammarlens}

\subsection*{Your turn}
\begin{itemize}
\raggedright
  \item \emph{Say it:} \emph{poruḷātāram}
  \item \emph{Say it:} \emph{poruḷātāram}, once more
  \item \emph{Say it:} \emph{maḻai peyyāviṭṭāl nāṭṭiṉ poruḷātāram pātikkappaṭum.} — If it does not rain, the country's economy will suffer.
  \item \emph{Recall:} say the Tamil for to restructure, then the Tamil for to strengthen, then say \emph{poruḷātāram} again
\end{itemize}

\begin{tcolorbox}[breakable,colback=teal!4,colframe=teal!35!black,title={Before you move on}]
What does \ta{பொருளாதாரம்} mean? (\textquotedblleft{}The economy\textquotedblright{}.) Say it once more.
\end{tcolorbox}
\section[\texorpdfstring{\ta{மானியம்} (māṉiyam)}{மானியம் (māṉiyam)}]{\ta{மானியம்} (māṉiyam) — a subsidy; a grant}

\label{lesson:TA-C391-maniyam}



\begin{quote}
Before the new one: say the Tamil for to strengthen, then the Tamil for the economy.
\end{quote}

\subsection*{You'll want to know: \ta{மானியம்}}
\textbf{\ta{மானியம்}} — \emph{māṉiyam} — \textquotedblleft{}a subsidy; a grant\textquotedblright{}.

\begin{quote}
\ta{அரசு} \ta{மானியம்} \ta{கொடுத்தால்} \ta{விவசாயிகள்} \ta{பயன்பெறுவார்கள்}.
\end{quote}

\begin{quote}
\emph{aracu māṉiyam koṭuttāl vivacāyikaḷ payaṉpeṟuvārkaḷ.}
\end{quote}

\begin{quote}
If the government gives a subsidy, the farmers will benefit.
\end{quote}

In older usage, land granted free of rent.

\begin{cousinweb}
\textbf{\ta{மானியம்}} (\emph{māṉiyam}) is listed in DEDR 4674 with Tamil \emph{maṇiyam}, the office of a village headman, and \emph{māṉiyam}, the rent-free land or payment he held: a grant from the state. Whether the word is native is not settled.
\end{cousinweb}

\begin{grammarlens}[title={What you've built}]
One more word for saying what happens if.
\end{grammarlens}

\subsection*{Your turn}
\begin{itemize}
\raggedright
  \item \emph{Say it:} \emph{māṉiyam}
  \item \emph{Say it:} \emph{māṉiyam}, once more
  \item \emph{Say it:} \emph{aracu māṉiyam koṭuttāl vivacāyikaḷ payaṉpeṟuvārkaḷ.} — If the government gives a subsidy, the farmers will benefit.
  \item \emph{Recall:} say the Tamil for to strengthen, then the Tamil for the economy, then say \emph{māṉiyam} again
\end{itemize}

\begin{tcolorbox}[breakable,colback=teal!4,colframe=teal!35!black,title={Before you move on}]
What does \ta{மானியம்} mean? (\textquotedblleft{}A subsidy; a grant\textquotedblright{}.) Say it once more.
\end{tcolorbox}
\section[\texorpdfstring{\ta{பற்றாக்குறை} (paṟṟākkuṟai)}{பற்றாக்குறை (paṟṟākkuṟai)}]{\ta{பற்றாக்குறை} (paṟṟākkuṟai) — a deficit; a shortage}

\label{lesson:TA-C391-parrakkurai}



\begin{quote}
Before the new one: say the Tamil for the economy, then the Tamil for a subsidy.
\end{quote}

\subsection*{You'll want to know: \ta{பற்றாக்குறை}}
\textbf{\ta{பற்றாக்குறை}} — \emph{paṟṟākkuṟai} — \textquotedblleft{}a deficit; a shortage\textquotedblright{}.

\begin{quote}
\ta{மழை} \ta{இல்லாவிட்டால்} \ta{தண்ணீர்ப்} \ta{பற்றாக்குறை} \ta{ஏற்படும்}.
\end{quote}

\begin{quote}
\emph{maḻai illāviṭṭāl taṇṇīrp paṟṟākkuṟai ēṟpaṭum.}
\end{quote}

\begin{quote}
If there is no rain, there will be a water shortage.
\end{quote}

\ta{பற்றா}, not enough + \ta{குறை}, a lack.

\begin{cousinweb}
\textbf{\ta{பற்றாக்குறை}} (\emph{paṟṟākkuṟai}) is native: \emph{paṟṟā}, \textquotedblleft{}not sufficing\textquotedblright{}, the negative of \emph{paṟṟu} (DEDR 4034), and \emph{kuṟai}, \textquotedblleft{}deficiency, lack\textquotedblright{} (DEDR 1851).
\end{cousinweb}

\begin{grammarlens}[title={What you've built}]
One more word for saying what happens if.
\end{grammarlens}

\subsection*{Your turn}
\begin{itemize}
\raggedright
  \item \emph{Say it:} \emph{paṟṟākkuṟai}
  \item \emph{Say it:} \emph{paṟṟākkuṟai}, once more
  \item \emph{Say it:} \emph{maḻai illāviṭṭāl taṇṇīrp paṟṟākkuṟai ēṟpaṭum.} — If there is no rain, there will be a water shortage.
  \item \emph{Recall:} say the Tamil for the economy, then the Tamil for a subsidy, then say \emph{paṟṟākkuṟai} again
\end{itemize}

\begin{tcolorbox}[breakable,colback=teal!4,colframe=teal!35!black,title={Before you move on}]
What does \ta{பற்றாக்குறை} mean? (\textquotedblleft{}A deficit; a shortage\textquotedblright{}.) Say it once more.
\end{tcolorbox}
\section[\texorpdfstring{\ta{விநியோகி} (viniyōki)}{விநியோகி (viniyōki)}]{\ta{விநியோகி} (viniyōki) — to distribute; to supply}

\label{lesson:TA-C391-viniyoki}



\begin{quote}
Before the new one: say the Tamil for a subsidy, then the Tamil for a deficit.
\end{quote}

\subsection*{You'll want to know: \ta{விநியோகி}}
\textbf{\ta{விநியோகி}} — \emph{viniyōki} — \textquotedblleft{}to distribute; to supply\textquotedblright{}.

\begin{quote}
\ta{அரசு} \ta{வெள்ளம்} \ta{வந்த} \ta{ஊர்களுக்கு} \ta{உணவு} \ta{விநியோகித்தது}.
\end{quote}

\begin{quote}
\emph{aracu veḷḷam vanta ūrkaḷukku uṇavu viniyōkittatu.}
\end{quote}

\begin{quote}
The government distributed food to the flooded towns.
\end{quote}

Past (I …): \ta{விநியோகித்தேன்}, \emph{viniyōkittēṉ}.

To whom: -\ta{க்கு}. The noun is \ta{விநியோகம்}, supply.

\begin{cousinweb}
\textbf{\ta{விநியோகி}} (\emph{viniyōki}) is a verb built on Sanskrit \emph{viniyoga}, \textquotedblleft{}apportionment, distribution, division\textquotedblright{} (Monier-Williams), with the loan-verb ending -\ta{இ}.
\end{cousinweb}

\begin{grammarlens}[title={What you've built}]
One more word for saying what happens if.
\end{grammarlens}

\subsection*{Your turn}
\begin{itemize}
\raggedright
  \item \emph{Say it:} \emph{viniyōki}
  \item \emph{Say it:} \emph{viniyōki}, once more
  \item \emph{Say it:} \emph{aracu veḷḷam vanta ūrkaḷukku uṇavu viniyōkittatu.} — The government distributed food to the flooded towns.
  \item \emph{Recall:} say the Tamil for a subsidy, then the Tamil for a deficit, then say \emph{viniyōki} again
\end{itemize}

\begin{tcolorbox}[breakable,colback=teal!4,colframe=teal!35!black,title={Before you move on}]
What does \ta{விநியோகி} mean? (\textquotedblleft{}To distribute; to supply\textquotedblright{}.) Say it once more.
\end{tcolorbox}
\section[\texorpdfstring{\ta{நிர்ணயி} (nirṇayi)}{நிர்ணயி (nirṇayi)}]{\ta{நிர்ணயி} (nirṇayi) — to fix; to set}

\label{lesson:TA-C391-nirnayi}



\begin{quote}
Before the new one: say the Tamil for a deficit, then the Tamil for to distribute.
\end{quote}

\subsection*{You'll want to know: \ta{நிர்ணயி}}
\textbf{\ta{நிர்ணயி}} — \emph{nirṇayi} — \textquotedblleft{}to fix; to set\textquotedblright{}.

\begin{quote}
\ta{அரசு} \ta{அரிசியின்} \ta{விலையை} \ta{நிர்ணயித்தது}.
\end{quote}

\begin{quote}
\emph{aracu ariciyiṉ vilaiyai nirṇayittatu.}
\end{quote}

\begin{quote}
The government fixed the price of rice.
\end{quote}

Past (I …): \ta{நிர்ணயித்தேன்}, \emph{nirṇayittēṉ}.

What is fixed takes -\ta{ஐ}: \ta{விலையை} \ta{நிர்ணயி}. The noun is \ta{நிர்ணயம்}.

\begin{cousinweb}
\textbf{\ta{நிர்ணயி}} (\emph{nirṇayi}) is a verb built on Sanskrit \emph{nirṇaya}, \textquotedblleft{}complete ascertainment, decision, settlement\textquotedblright{} (Monier-Williams), with the loan-verb ending -\ta{இ}.
\end{cousinweb}

\begin{grammarlens}[title={What you've built}]
One more word for saying what happens if.
\end{grammarlens}

\subsection*{Your turn}
\begin{itemize}
\raggedright
  \item \emph{Say it:} \emph{nirṇayi}
  \item \emph{Say it:} \emph{nirṇayi}, once more
  \item \emph{Say it:} \emph{aracu ariciyiṉ vilaiyai nirṇayittatu.} — The government fixed the price of rice.
  \item \emph{Recall:} say the Tamil for a deficit, then the Tamil for to distribute, then say \emph{nirṇayi} again
\end{itemize}

\begin{tcolorbox}[breakable,colback=teal!4,colframe=teal!35!black,title={Before you move on}]
What does \ta{நிர்ணயி} mean? (\textquotedblleft{}To fix; to set\textquotedblright{}.) Say it once more.
\end{tcolorbox}
